\documentclass[11pt,a4paper]{article}
\usepackage[T1]{fontenc}
\usepackage{lmodern}
\usepackage[margin=22mm,headheight=15pt]{geometry}
\usepackage{amsmath,amssymb,mathtools,microtype}
\usepackage{xcolor,enumitem,fancyhdr,booktabs,needspace}
\usepackage[hidelinks]{hyperref}
\definecolor{examblue}{HTML}{17365D}
\pagestyle{fancy}
\fancyhf{}
\lhead{\textsc{Machine Learning 1}}
\rhead{\textsc{Mock 2 --- B: Answer Key}}
\cfoot{\thepage}
\setlength{\parindent}{0pt}
\setlength{\parskip}{0.4em}
\setlist[itemize]{leftmargin=1.5em,itemsep=0.1em}
\newcommand{\R}{\mathbb R}
\newcommand{\E}{\mathbb E}
\DeclareMathOperator{\Var}{Var}
\DeclareMathOperator{\ReLU}{ReLU}
\newcommand{\solution}[3]{\Needspace{6\baselineskip}\subsection*{\color{examblue}#1\quad #2\hfill\normalsize[#3 p]}}
\newcommand{\selfcredit}[1]{\par\textit{Suggested credit: #1}\par}

\begin{document}
\begin{center}
{\Huge\bfseries\color{examblue}Mock Midterm 2 --- B}\par
\vspace{0.5em}
{\LARGE Worked solutions and marking guide}\par
\vspace{0.7em}
{\large 18 points: Exercise 1 (8), Exercise 2 (10)}
\end{center}
These are full solutions, not only final answers. The marking allocations are
for self-assessment, not an official course rubric. Any equivalent network
construction or correct derivation can receive full credit. Give partial credit
for a correct method even if a later arithmetic error propagates.

\section*{Exercise 1: A nonlinear Bayes rule and a small neural network}
\solution{1a}{Posterior log odds and a bent boundary}{2}
Bayes' rule cancels the evidence:
\[
a(x)=\log\frac{\pi_1}{\pi_2}
+\log\frac{p(x_1\mid C_1)}{p(x_1\mid C_2)}
+\log\frac{p(x_2\mid C_1)}{p(x_2\mid C_2)}.
\]
The prior term is $-\log2$. The Laplace ratio contributes
\[
\log2-|x_1|+\tfrac12|x_1|=\log2-\tfrac12|x_1|.
\]
The Gaussian normalizers cancel, and their log ratio is
\[
-\tfrac12\left[(x_2-\tfrac12)^2-(x_2+\tfrac12)^2\right]=x_2.
\]
All constants therefore cancel:
\[
\boxed{a(x)=x_2-\tfrac12|x_1|,\qquad
p(C_1\mid x)=\sigma(x_2-\tfrac12|x_1|)}.
\]
The strict decision region is
\[
\boxed{R_1=\{x:x_2>\tfrac12|x_1|\}}.
\]
The boundary is a V with vertex at the origin: the ray $x_2=x_1/2$ for
$x_1\geq0$ and the ray $x_2=-x_1/2$ for $x_1\leq0$. The region is above
both rays. On the boundary, both posteriors are $1/2$.

Naive Bayes says that the log odds are a sum of one-coordinate functions.
It does \emph{not} require each function to be linear in its coordinate.
The Laplace log ratio introduces an absolute value, giving a piecewise-linear
rather than a globally affine decision rule.
\selfcredit{1 p for the simplified posterior/log odds; 0.5 p for the region and boundary; 0.5 p for the distinction between conditional independence and linearity.}

\solution{1b}{Correct labels, nonlinear features, and nonexistence of a finite MLE}{2}
Let an affine logit be $g(x)=w_0+w_1x_1+w_2x_2$, with $g>0$ for $C_1$.
Because $(0,1)$ is the midpoint of the other two points,
\[
g(0,1)=\tfrac12\big[g(-4,1)+g(4,1)\big].
\]
If both outer points have strictly negative logits, their midpoint does too.
This contradicts the required positive logit at the middle point. No affine
classifier can strictly classify the three points.

One valid fixed feature map and weight vector are
\[
\boxed{\phi(x)=(1,|x_1|,x_2)^\top,\qquad
w=(0,-\tfrac12,1)^\top}.
\]
They give $w^\top\phi(x)=a(x)$ for every input, not only the three points.
The three logits are $+1,-1,-1$, so the transformed data are strictly
linearly separable. Multiplying $w$ by $q>0$ preserves all predicted labels
and increases their correct-class probabilities. The negative log-likelihood is
\[
L(qw)=3\log(1+e^{-q})\longrightarrow0\quad\text{as }q\to\infty.
\]
At any finite weight vector, each sigmoid probability lies strictly between
zero and one, so every correctly labeled observation has a strictly positive
loss. The likelihood supremum is 1 but is not attained at finite weights:
\[
\boxed{\text{There is no finite unregularized MLE for this separable dataset.}}
\]
Finite weights can give perfect classification accuracy without assigning
probability 1 to the observed labels. A positive definite Gaussian weight prior
would prevent this unlimited scaling for logistic regression with fixed features.
\selfcredit{0.75 p for the affine impossibility proof; 0.5 p for exact nonlinear features; 0.75 p for the finite-MLE distinction and scaling argument.}

\solution{1c}{An exact two-unit construction with a direct connection}{1}
Use
\[
\boxed{U=\begin{pmatrix}1&0\\-1&0\end{pmatrix},\quad
b=\begin{pmatrix}0\\0\end{pmatrix},\quad
c=\begin{pmatrix}-1/2\\-1/2\end{pmatrix},\quad
v=\begin{pmatrix}0\\1\end{pmatrix},\quad d=0}.
\]
Then $h_1=\ReLU(x_1)$ and $h_2=\ReLU(-x_1)$, and
\[
c^\top h+v^\top x+d
=-\tfrac12[\ReLU(x_1)+\ReLU(-x_1)]+x_2
=x_2-\tfrac12|x_1|.
\]
The hidden units form the absolute-value term. The direct connection supplies
$x_2$ without using extra hidden units. All parameters are finite. Hidden-unit
permutations and positive rescalings of ReLU units give equivalent constructions.
\selfcredit{0.75 p for a valid exact construction; 0.25 p for identifying the two paths.}

\solution{1d}{Backpropagation through both paths}{2}
First differentiate the binary cross entropy and sigmoid:
\[
\frac{\partial L}{\partial y}=-\frac ty+\frac{1-t}{1-y},\qquad
\frac{dy}{da}=y(1-y),\qquad
\boxed{\delta:=\frac{\partial L}{\partial a}=y-t}.
\]
Let $I_i=\mathbf1\{r_i>0\}$ and let $D_h=\operatorname{diag}(I_1,I_2)$.
From $r_i=\sum_jU_{ij}x_j+b_i$ and
$a=\sum_ic_ih_i+\sum_jv_jx_j+d$, component derivatives are
\begin{align*}
\frac{\partial L}{\partial U_{ij}}&=\delta c_i I_i x_j,&
\frac{\partial L}{\partial b_i}&=\delta c_i I_i,\\
\frac{\partial L}{\partial c_i}&=\delta h_i,&
\frac{\partial L}{\partial v_j}&=\delta x_j,&
\frac{\partial L}{\partial d}&=\delta,\\
\frac{\partial L}{\partial x_j}&=
\delta\left(\sum_ic_iI_iU_{ij}+v_j\right).
\end{align*}
Assembled gradients and the numerator-layout input derivative are
\begin{align*}
\nabla_U L&=\delta(D_hc)x^\top\in\R^{2\times2},&
\nabla_b L&=\delta D_hc\in\R^2,\\
\nabla_c L&=\delta h\in\R^2,&
\nabla_v L&=\delta x\in\R^2,\\
\boxed{\frac{dL}{dx}}&\boxed{=\delta(c^\top D_hU+v^\top)\in\R^{1\times2}}.
\end{align*}
The input derivative is the \emph{sum} of the hidden-layer path and the
direct path; omitting $v^\top$ gives the wrong result.

For the specified numerical check,
\[
r=(1,-2)^\top,\quad h=(1,0)^\top,\quad a=0,\quad
y=\tfrac12,\quad \delta=-\tfrac12,\quad D_hc=(1,0)^\top.
\]
Therefore
\[
\boxed{\nabla_U L=\begin{pmatrix}-1&0\\0&0\end{pmatrix},\qquad
\frac{dL}{dx}=(-\tfrac12,-1)}.
\]
For reference, the other gradients are
$\nabla_bL=(-\tfrac12,0)^\top$, $\nabla_cL=(-\tfrac12,0)^\top$,
$\nabla_vL=(-1,0)^\top$, and $\partial L/\partial d=-\tfrac12$.
\selfcredit{0.25 p for the logit derivative; 1.25 p for component and assembled derivatives, including both input paths; 0.5 p for the two numerical results.}

\solution{1e}{Zero initialization cannot learn the hidden features}{1}
At initialization $U=b=c=0$, so $r=h=0$ for every training input. Under
$\ReLU'(0)=0$, $D_h=0$. The formulas in 1d give
\[
\nabla_U L=0,\qquad\nabla_b L=0,\qquad\nabla_c L=\delta h=0
\]
for every observation, regardless of the current $v,d$.
An isotropic Gaussian prior adds gradients proportional to the parameters,
which are also zero for $U,b,c$. A gradient step therefore leaves all three
at zero. This repeats at every subsequent step, proving the claim by induction.

The resulting predictor is always $\sigma(v^\top x+d)$, an affine-logit
classifier. By 1b it cannot strictly classify all three points correctly.
Weight decay cannot move a parameter away from zero when both its data
gradient and its penalty gradient are zero. A nonzero, appropriately scaled
hidden initialization avoids this particular frozen state; this is an
optimization failure, not a lack of representational capacity of the full network.
\selfcredit{0.5 p for the gradient/induction argument; 0.5 p for the capacity reduction and why the prior cannot revive the layer.}

\newpage
\section*{Exercise 2: Soft routing among regression experts}
\solution{2a}{Joint, marginal, and posterior expert probabilities}{2}
Write $m_{nk}=\theta_k^\top\phi_n$ and
$f_{nk}=\mathcal N(t_n\mid m_{nk},s^2)$. With one-hot expert indicators,
\[
\boxed{p(t,Z\mid X,\Theta,V,s^2)=
\prod_{n=1}^N\prod_{k=1}^K[\pi_{nk}f_{nk}]^{z_{nk}}}.
\]
For each observation exactly one expert is active. Summing over its possible
values gives
\[
\boxed{p(t\mid X,\Theta,V,s^2)=\prod_n\left(\sum_k\pi_{nk}f_{nk}\right)},
\qquad
\boxed{\ell=\sum_n\log\left(\sum_k\pi_{nk}f_{nk}\right)}.
\]
By Bayes' rule,
\[
\boxed{r_{nk}=\frac{\pi_{nk}f_{nk}}{\sum_j\pi_{nj}f_{nj}}},\qquad
\sum_kr_{nk}=1.
\]
The experts are alternative explanations for one target, not independent
observations of that target. Their densities must be added with routing
weights. Multiplying them neither marginalizes a categorical latent variable
nor generally gives the stated normalized mixture density.
\selfcredit{0.5 p each for complete likelihood, observed likelihood/log-likelihood, responsibilities, and the sum-versus-product explanation.}

\solution{2b}{Observed-likelihood gradients and an unidentifiable gate}{2}
Let $p_n=\sum_j\pi_{nj}f_{nj}$. For an expert parameter component,
\begin{align*}
\frac{\partial\ell}{\partial\theta_{kr}}
&=\sum_n\frac{\pi_{nk}f_{nk}}{p_n}
\frac{\partial\log f_{nk}}{\partial\theta_{kr}}\\
&=\sum_nr_{nk}\frac{t_n-\theta_k^\top\phi_n}{s^2}\phi_{nr}.
\end{align*}
Thus the column gradient is
\[
\boxed{\nabla_{\theta_k}\ell
=\frac1{s^2}\sum_nr_{nk}(t_n-\theta_k^\top\phi_n)\phi_n\in\R^M}.
\]
For the routing parameters, differentiate the softmax in components:
\[
\frac{\partial\pi_{nj}}{\partial v_{kr}}
=\pi_{nj}(\delta_{jk}-\pi_{nk})\phi_{nr}.
\]
Consequently
\begin{align*}
\frac{\partial\ell}{\partial v_{kr}}
&=\sum_n\frac1{p_n}\sum_jf_{nj}\pi_{nj}
(\delta_{jk}-\pi_{nk})\phi_{nr}\\
&=\sum_n\left(r_{nk}-\pi_{nk}\sum_jr_{nj}\right)\phi_{nr}
=\sum_n(r_{nk}-\pi_{nk})\phi_{nr},
\end{align*}
so
\[
\boxed{\nabla_{v_k}\ell=\sum_n(r_{nk}-\pi_{nk})\phi_n\in\R^M}.
\]
These formulas come directly from differentiating the log of a sum. Although
the responsibilities depend on current parameters, they appear as the result
of the derivative; there is no additional $\partial r/\partial\theta$ term in
these final formulas.

If all expert densities are the same, say $f_{nk}=f_n$, then
$p_n=f_n\sum_k\pi_{nk}=f_n$ and $r_{nk}=\pi_{nk}$. Hence every routing
gradient is zero and the likelihood is independent of $V$. The data provide
no information about routing between identical experts. This is a stronger
ambiguity than the usual softmax common-shift invariance.
\selfcredit{0.75 p for each gradient with a component derivation; 0.5 p for the identical-expert conclusion.}

\solution{2c}{EM updates and the condition for monotonicity}{3}
Initialize finite expert and routing parameters and $s^2>0$, preferably using
different initial experts. At iteration $j$:

\textbf{E-step.} Use the old parameters to calculate
\[
r_{nk}^{(j)}=
\frac{\pi_{nk}^{(j)}\mathcal N(t_n\mid(\theta_k^{(j)})^\top\phi_n,(s^2)^{(j)})}
{\sum_\ell\pi_{n\ell}^{(j)}\mathcal N(t_n\mid(\theta_\ell^{(j)})^\top\phi_n,(s^2)^{(j)})}.
\]
\textbf{M-step or generalized M-step.} Hold these responsibilities fixed in
\[
Q(\Theta,V,s^2)=\sum_{n,k}r_{nk}^{(j)}\left[
\log\pi_{nk}(V)-\frac12\log(2\pi s^2)
-\frac{(t_n-\theta_k^\top\phi_n)^2}{2s^2}\right].
\]
Do not recompute or differentiate the responsibilities inside this objective.
The old posterior defines the expectation; the new parameters are its
optimization variables.

Let $\Phi\in\R^{N\times M}$ have rows $\phi_n^\top$,
$t\in\R^N$, and $D_k=\operatorname{diag}(r_{1k}^{(j)},\ldots,r_{Nk}^{(j)})$.
Differentiation of $Q$ with respect to an expert component gives
\[
\frac{\partial Q}{\partial\theta_{kr}}
=\frac1{s^2}\sum_nr_{nk}^{(j)}(t_n-\theta_k^\top\phi_n)\phi_{nr}.
\]
Setting this to zero yields the weighted normal equations and update
\[
\Phi^\top D_k\Phi\,\theta_k=\Phi^\top D_kt,
\qquad
\boxed{\theta_k^{(j+1)}=(\Phi^\top D_k\Phi)^{-1}\Phi^\top D_kt}.
\]
The inverse exists by assumption. This weighted least-squares update is
independent of the shared variance. The factor $1/s^2$ cancels, but the
responsibilities still depended on the old variance in the E-step.

Set
$A=\sum_{n,k}r_{nk}^{(j)}[t_n-(\theta_k^{(j+1)})^\top\phi_n]^2$.
Because $\sum_{n,k}r_{nk}^{(j)}=N$, the variance-dependent part of $Q$ is
$-\tfrac N2\log s^2-A/(2s^2)$. Treating $s^2$ as the optimization variable,
\[
\frac{\partial Q}{\partial s^2}=-\frac N{2s^2}+\frac A{2(s^2)^2}=0,
\qquad
\boxed{(s^2)^{(j+1)}=\frac1N\sum_{n,k}r_{nk}^{(j)}
[t_n-(\theta_k^{(j+1)})^\top\phi_n]^2}.
\]
The denominator is $N$, not $NK$, since the responsibilities of each
observation sum to one. Positivity of $A$ gives an interior variance maximum.

The routing objective is
$Q_V=\sum_{n,k}r_{nk}^{(j)}\log\pi_{nk}(V)$, a soft-target multinomial
logistic log-likelihood. Its gradient is
\[
\nabla_{v_k}Q_V=\sum_n(r_{nk}^{(j)}-\pi_{nk}(V))\phi_n.
\]
One valid generalized-EM update starts at the old $V$ and uses
\[
v_k^{\mathrm{new}}=v_k^{\mathrm{old}}
+\eta\sum_n(r_{nk}^{(j)}-\pi_{nk}(V^{\mathrm{old}}))\phi_n,
\]
with backtracking on $\eta>0$ until $Q_V(V^{\mathrm{new}})\geq Q_V(V^{\mathrm{old}})$.
Further steps may be taken with the \emph{same fixed} responsibilities.
A gradient direction alone is not a guarantee: an excessive step can reduce
$Q_V$. Increasing the frozen-responsibility $Q$ gives the standard EM lower-bound
argument for a nondecreasing observed-data likelihood. Exact expert/variance
maximization and a verified nondecreasing routing block suffice.

After the updates, evaluate the new observed-data log-likelihood. For example,
stop if its nonnegative relative improvement is less than $\varepsilon$,
\[
\frac{\ell^{(j+1)}-\ell^{(j)}}{1+|\ell^{(j)}|}<\varepsilon,
\]
subject to a maximum iteration count and a numerical check that the likelihood
did not decrease beyond round-off tolerance. A parameter-change or gradient
check can also be used. Multiple initializations reduce sensitivity to local
solutions. EM/GEM does not guarantee a global optimum of this nonconvex
observed-likelihood problem; monotonicity is not global optimality.
\selfcredit{0.5 p for the E-step and frozen-responsibility $Q$; 0.75 p each for expert and variance updates; 1 p for valid routing ascent, a convergence check, and the qualified monotonicity/global-optimum statements.}

\solution{2d}{Routing and responsibility can select different experts}{1}
Before seeing the target, hard routing chooses $\boxed{\text{expert 1}}$
because $4/5>1/5$. After seeing $t=2$, the two Gaussian densities are
\[
f_1=\frac1{\sqrt{2\pi}}e^{-9/2},\qquad f_2=\frac1{\sqrt{2\pi}}.
\]
Thus
\[
\boxed{r_1=\frac{4e^{-9/2}}{1+4e^{-9/2}}\approx0.04255,\qquad
r_2=\frac1{1+4e^{-9/2}}\approx0.95745}.
\]
The highest-responsibility choice is now $\boxed{\text{expert 2}}$.
Routing uses only the input. Responsibility also uses how well each expert
explains the observed target; expert 2's much larger density outweighs its
smaller routing probability.
\selfcredit{0.25 p for hard routing; 0.5 p for responsibilities and posterior selection; 0.25 p for the explanation.}

\solution{2e}{Predictive uncertainty and squared-error prediction}{2}
At a new input $x$, let $\pi_k=\pi_k(x)$ and
$m_k=\theta_k^\top\phi(x)$. By conditioning on the latent expert,
\[
\boxed{\bar m(x)=\E[T\mid x]=\sum_k\pi_km_k}.
\]
Each expert has second moment $s^2+m_k^2$. Therefore
\begin{align*}
\Var(T\mid x)
&=\sum_k\pi_k(s^2+m_k^2)-\bar m^2\\
&=\boxed{s^2+\sum_k\pi_k(m_k-\bar m)^2}.
\end{align*}
The first term is within-expert noise. The second is between-expert
uncertainty from not knowing which expert generates the target. With
expert-specific variances the first term would instead be
$\sum_k\pi_ks_k^2$; here it is simply $s^2$.

For any scalar prediction $q(x)$,
\[
\E[(T-q)^2\mid x]
=\Var(T\mid x)+(\bar m-q)^2,
\]
because the cross term involving $T-\bar m$ has zero conditional expectation.
Thus the optimal squared-error prediction is $\boxed{q(x)=\bar m(x)}$,
not necessarily the mean of the most likely expert.

In the numerical case,
\[
\boxed{\bar m=\tfrac45(-1)+\tfrac15(2)=-\tfrac25=-0.4},
\]
\[
\boxed{\Var(T\mid x)=1+\tfrac45(-\tfrac35)^2+\tfrac15(\tfrac{12}{5})^2
=\frac{61}{25}=2.44}.
\]
Responsibilities $p(z=k\mid x,t)$ require the very target that is being
predicted. Substituting the values from 2d would use target information that
is unavailable before observing $T$. For an unobserved target, marginalize
the expert using $p(z=k\mid x)$, the routing probabilities.
\selfcredit{0.5 p for mean and variance; 0.5 p for the squared-error argument; 0.5 p for the numerical values; 0.5 p for distinguishing routing from target-conditioned responsibilities.}

\section*{Scope and source alignment}
The format follows \path{exams/homework-based-exam-1.pdf}: two extended
exercises, two hours, 18 points, derivations, and follow-up insights. The
actual second midterm's question count and topic selection are unknown.
\begin{itemize}
  \item Exercise 1 develops HW 2 Exercises 1--2, practice sets 3--4, and lectures 6--8. New elements are Laplace/Gaussian class conditionals, a V-shaped Bayes boundary, a new three-point obstruction, exact ReLU features with a direct connection, derivatives through two input paths, and frozen hidden initialization.
  \item Exercise 2 develops HW 2 Exercise 6, practice set 5, and lecture 9. It replaces the exponential experts with Gaussian regression experts and adds weighted normal equations, a learned shared variance, a verified generalized-EM routing step, indistinguishable experts, and predictive variance.
  \item Together with Mock A, the pair covers all six HW 2 themes: Naive Bayes, classifier capacity, regularized logistic regression, PCA, generalized Gaussian latent coordinates in PPCA, and mixtures of experts.
\end{itemize}
The foundational operations necessarily recur, but no original homework or
practice-set problem, dataset, or complete derivation sequence is reproduced.
The questions are intentionally harder than the first midterm; the point total
does not imply equal workload.
\end{document}
